\booksection{第十一章附：專業影像技術知識}

\chapterintro{想要真正理解手機攝影同影片嘅奧秘，你需要了解一啲專業嘅影像技術知識。呢啲知識唔淨係可以幫你更好咁使用手機相機，仲可以等你喺選擇設備嗰陣更有判斷力。}

\booksubsection{一、分辨率、幀率與碼率：影片質量嘅三要素}

\subsubsection*{分辨率（Resolution）：}

分辨率決定咗畫面嘅清晰度，用像素數表示。

\begin{itemize}
  \item \textbf{720P (HD)}：1280×720像素，入門級清晰度。
  \item \textbf{1080P (Full HD)}：1920×1080像素，目前最主流嘅分辨率。
  \item \textbf{2K (QHD)}：2560×1440像素，比1080P更清晰。
  \item \textbf{4K (UHD)}：3840×2160像素，超高清，細節豐富。
  \item \textbf{8K}：7680×4320像素，未來標準，目前設備支援有限。
\end{itemize}

分辨率越高，畫面越清晰，但檔案都越大，對存儲同處理能力要求越高。

\subsubsection*{幀率（Frame Rate）：}

幀率決定咗畫面嘅流暢度，用每秒幀數（fps）表示。

\begin{itemize}
  \item \textbf{24fps}：電影標準幀率，有獨特嘅「電影感」。
  \item \textbf{30fps}：電視同網絡影片常用，流暢自然。
  \item \textbf{60fps}：高幀率，適合運動場景同遊戲錄製，非常流暢。
  \item \textbf{120fps/240fps}：超高幀率，可用於慢動作影片。
\end{itemize}

幀率越高，畫面越流暢，慢動作效果越好，但檔案都越大。

\subsubsection*{碼率（Bitrate）：}

碼率決定咗影片嘅數據密度，用每秒數據量（Mbps）表示。

\begin{itemize}
  \item 碼率越高，畫質越好，細節保留越多，但檔案都越大。
  \item 碼率過低會導致畫面出現馬賽克、色塊等壓縮痕跡。
  \item 專業影片拍攝通常需要高碼率以保留更多後期處理空間。
\end{itemize}

\booksubsection{二、影片編碼格式：H.264、H.265與ProRes}

\subsubsection*{H.264 (AVC)：}

\begin{itemize}
  \item \textbf{特點}：最廣泛使用嘅編碼格式，兼容性極好。
  \item \textbf{優點}：幾乎所有設備都可以播放，壓縮效率適中。
  \item \textbf{缺點}：相比新格式，壓縮效率較低，檔案較大。
  \item \textbf{適用場景}：網絡影片分享、普通用戶日常拍攝。
\end{itemize}

\subsubsection*{H.265 (HEVC)：}

\begin{itemize}
  \item \textbf{特點}：H.264嘅升級版，壓縮效率提升約50\%。
  \item \textbf{優點}：相同畫質下檔案更細，支援更高分辨率同幀率。
  \item \textbf{缺點}：編碼同解碼需要更強嘅處理器，部分老設備唔兼容。
  \item \textbf{適用場景}：4K影片拍攝、存儲空間有限嘅場景。
\end{itemize}

\subsubsection*{Apple ProRes：}

\begin{itemize}
  \item \textbf{特點}：蘋果開發嘅專業影片編碼，幾乎無損壓縮。
  \item \textbf{優點}：畫質極高，後期處理空間大，色彩準確。
  \item \textbf{缺點}：檔案極大，需要大量存儲空間，處理需要專業設備。
  \item \textbf{適用場景}：專業影片製作、電影級拍攝。
\end{itemize}

\subsubsection*{ProRes RAW：}

\begin{itemize}
  \item \textbf{特點}：結合咗ProRes同RAW格式嘅優勢。
  \item \textbf{RAW係乜嘢}：未經處理嘅原始傳感器數據，保留最大資訊量。
  \item \textbf{優點}：後期調整空間極大，可以改變曝光、白平衡等參數而唔損失畫質。
  \item \textbf{缺點}：檔案極其巨大，需要專業後期軟件處理。
  \item \textbf{適用場景}：高端專業製作。iPhone 15 Pro系列及以後支援拍攝Apple ProRes影片（最高4K 60fps），iPhone 17 Pro進一步支援ProRes RAW直接機內錄製。
\end{itemize}

\booksubsection{三、HDR：高動態範圍技術}

HDR（High Dynamic Range）技術可以同時保留畫面中最亮同最暗部分嘅細節。

\begin{itemize}
  \item \textbf{HDR10}：最基礎嘅HDR標準，兼容性好。
  \item \textbf{HDR10+}：三星推動嘅標準，支援動態元數據。
  \item \textbf{Dolby Vision}：杜比開發嘅頂級HDR標準，畫質最佳，但需要授權。
  \item \textbf{HLG}：廣播級HDR標準，兼容普通顯示器。
\end{itemize}

拍攝HDR影片嗰陣，需要支援HDR嘅顯示器先可以睇到完整效果。

\booksubsection{四、色彩空間與色深}

\subsubsection*{色彩空間：}

\begin{itemize}
  \item \textbf{sRGB}：標準色彩空間，兼容性最好。
  \item \textbf{Display P3}：蘋果推廣嘅廣色域標準，色彩更豐富。
  \item \textbf{Rec.2020}：HDR影片嘅標準色彩空間，色域最廣。
\end{itemize}

\subsubsection*{色深：}

\begin{itemize}
  \item \textbf{8-bit}：1670萬色，普通影片標準。
  \item \textbf{10-bit}：10.7億色，專業影片標準，色彩過渡更平滑。
  \item \textbf{12-bit}：680億色，高端專業標準。
\end{itemize}

\begin{tcbinfo}{實用建議}
\begin{itemize}
  \item 日常拍攝：1080P/4K + 30fps + H.265編碼，平衡畫質同檔案大小。
  \item 運動拍攝：1080P/4K + 60fps，保證流暢度。
  \item 專業創作：4K + ProRes（如果支援），保留最大後期空間。
  \item 存儲建議：高碼率影片好佔空間，建議定期備份到電腦或雲盤。
\end{itemize}

\end{tcbinfo}

\begin{tcbpractice}{去探索一下：你嘅手機影片設定}
\begin{enumerate}
  \item 打開相機App，切換到影片模式。
  \item 查看分辨率同幀率選項，嘗試唔同嘅組合。
  \item 如果支援，查看有冇HDR、ProRes等高級選項。
  \item 拍攝幾段唔同設定嘅影片，對比畫質同檔案大小。
\end{enumerate}

\end{tcbpractice}

\booksubsection{五、RAW格式與HEIF/JPEG對比}

\subsubsection*{乜嘢係RAW格式：}

RAW格式係相機傳感器記錄嘅「原始數據」——未經任何壓縮同處理嘅純原始資訊。佢就好似生鮮食材，需要後期「烹飪」先可以成品。手機因為存儲同算力有限，預設影出嚟嘅都已經係「烹飪好」嘅JPEG或HEIF，RAW需要專門開啟。

\subsubsection*{RAW vs HEIF vs JPEG對比：}

\begin{itemize}
  \item \textbf{JPEG}：最老嘅格式，兼容性最好（任何設備都可以打開）。8-bit色深，有損壓縮，檔案最小（一張相2-5MB）。後期空間有限，強行提亮暗部容易出現色塊。
  \item \textbf{HEIF/HEIC}：蘋果主推嘅新格式，同樣畫質下檔案淨係得JPEG嘅一半（一張相1-3MB）。支援10-bit色深同透明通道。缺點係部分老設備同Windows需要安裝插件先可以打開。
  \item \textbf{RAW}：未經壓縮嘅原始數據，檔案最大（一張相20-100MB）。色深通常12-16bit，後期空間最大。缺點係必須用專業軟件（Lightroom、Snapseed等）先可以打開同處理，無法直接分享。
\end{itemize}

\subsubsection*{點解專業攝影師要影RAW：}

\begin{itemize}
  \item \textbf{白平衡可後期無損調整}：JPEG影錯咗白平衡（比如白熾燈下偏黃），後期糾正會損失畫質；RAW可以重新設定色温，完全無損。
  \item \textbf{曝光容差大}：JPEG欠曝2檔嘅暗部基本係死黑一片，提亮後全係噪點；RAW可以由欠曝3-4檔嘅暗部「救」返可用細節。過曝同理。
  \item \textbf{暗部同高光細節更豐富}：RAW嘅動態範圍遠大過JPEG，亮部不過曝、暗部不死黑嘅餘地大得多。
  \item \textbf{鋭化、降噪、鏡頭矯正都喺後期可控}：相機嘅機內處理係「一錘子買賣」，RAW可以反覆調整，每次都由原始數據出發。
\end{itemize}

\subsubsection*{手機點樣開啟RAW：}

\begin{itemize}
  \item \textbf{iPhone ProRAW}：僅iPhone 12 Pro及以上Pro機型支援。打開「設定 相機 格式」中開啟「Apple ProRAW」。然後喺相機介面頂部會出現一個「RAW」按鈕，點亮即可。檔案約25MB/張，建議有需要嗰陣再開。
  \item \textbf{安卓Motion Raw / Ultra RAW}：各品牌叫法唔同。小米/OPPO/vivo旗艦機型喺相機「專業模式」或「更多」入面可以揾到RAW開關；華為喺「專業模式」中開啟RAW；三星Galaxy S系列喺「專業」模式入面有RAW選項。檔案大小通常喺20-50MB。
  \item \textbf{通用建議}：日常隨手影冇必要開RAW，檔案太大佔空間；淨係喺重要場景（風光、人像、需要後期）嗰陣先開啟。
\end{itemize}

\booksubsection{六、計算攝影原理}

現代手機攝影已經離不開「計算攝影」——亦即係靠算法而唔淨係鏡頭同傳感器嚟生成相片。理解呢幾個核心原理，可以等你更明白手機喺撳快門嗰一刻做咗乜嘢。

\subsubsection*{多幀合成：}

手機撳一次快門，其實影咗多張相。算法將多幀對齊、加權融合，得到一張高質量成片。

\begin{itemize}
  \item \textbf{夜景模式}：1-3秒內連拍多張欠曝相，對齊後疊加——噪點互相抵消、亮度累加。結果就係暗光都可以影出乾淨嘅畫面。
  \item \textbf{HDR（高動態範圍）}：分別拍攝欠曝、正常、過曝三張（或多張），合成一張亮部同暗部都有細節嘅相。逆光場景尤其有效。
  \item \textbf{Deep Fusion（蘋果）/Photonic Engine（光子引擎）}：針對中低光場景，逐像素分析多幀並融合，重點提升紋理同細節，冷衫、毛髮呢類細節會明顯更清晰。
\end{itemize}

\subsubsection*{語義分割：}

AI唔淨係將成張圖一刀切處理，而係先「睇懂」畫面入面有乜嘢——天空、人臉、草地、建築、衫——然後對每個區域分別應用唔同嘅處理策略。

\begin{itemize}
  \item \textbf{天空}：單獨增強藍色、壓低亮度，避免過曝。
  \item \textbf{人臉}：識別膚色區域，單獨提亮、平滑膚色，唔動五官同頭髮細節。
  \item \textbf{草地/樹葉}：增強綠色飽和度，提升細節鋭化。
  \item \textbf{背景虛化}：人像模式下，AI識別主體輪廓，背景模糊、主體清晰。
\end{itemize}

\subsubsection*{深度學習嘅應用：}

\begin{itemize}
  \item \textbf{AI場景識別}：相機自動識別「美食」「寵物」「夜景」「文檔」等場景，並對應優化參數。比如美食模式會增強暖色調同飽和度，等食物睇落更誘人。
  \item \textbf{AI人像虛化}：手機鏡頭物理上無法做到單反嗰種淺景深，靠深度學習模型預測主體邊緣，再對背景做模糊。最新算法已經可以正確處理頭髮絲、玻璃邊緣呢啲細節。
  \item \textbf{AI消除/AI擴圖}：選中畫面中嘅路人或雜物，AI基於上下文「補全」被遮擋嘅部分。呢個係生成式AI喺影像領域嘅應用。
\end{itemize}

\subsubsection*{計算攝影 vs 光學攝影：}

\begin{itemize}
  \item \textbf{計算攝影嘅優勢}：細體積都可以得到好畫質，暗光、逆光等極端場景都可以出片，普通用戶唔需要懂攝影知識就可以影出「好睇」嘅相。
  \item \textbf{計算攝影嘅劣勢}：畫面有時「過度處理」——鋭化過度、色彩失真、天空顏色唔自然；多幀合成喺主體移動嗰陣會留低重影；細節係「算出嚟」嘅而唔係真實記錄嘅。
  \item \textbf{光學攝影嘅優勢}：大底傳感器+大光圈鏡頭能真實捕捉光線，畫質自然，後期空間大；唔會因為算法「腦補」產生偽影。
  \item \textbf{趨勢}：現代旗艦手機係「光學+計算」結合——用大底傳感器獲取真實資訊，用算法優化呈現。兩者並不對立。
\end{itemize}

\booksubsection{七、鏡頭類型與焦段}

現代旗艦手機通常配三到四顆鏡頭，每顆對應唔同焦段。理解焦段，可以等你喺唔同場景下揀啱鏡頭。

\subsubsection*{常見鏡頭與等效焦距：}

\begin{itemize}
  \item \textbf{超廣角（13-24mm等效）}：視野最廣，適合風光、建築、室內空間。能影到眼前所見嘅大場面，但邊緣會有畸變（直線變彎），影人像嗰陣人臉會被拉伸。適合「大場景」，唔適合「特寫」。
  \item \textbf{主攝（24-28mm等效）}：日常最常用嘅鏡頭，畫質最好、光圈最大。兼顧視野同畫質，影人影景都可以用。多數手機嘅「1x」就係主攝。
  \item \textbf{長焦（70-200mm等效）}：用於影遠景、人像特寫。長焦能壓縮空間感（遠處嘅山同近處嘅人睇落更近），影人像能獲得更自然嘅五官比例同背景虛化。
  \item \textbf{潛望式長焦（200mm+等效）}：通過反射式光路設計，喺手機厚度內實現更長焦距。能喺好遠距離影清楚細節，適合演唱會、體育賽事、野生動物。旗艦機型先有。
\end{itemize}

\subsubsection*{傳感器尺寸與像素大小嘅權衡：}

傳感器係「底」，決定能接收幾多光線。行業入面講「底大一級壓死人」——同樣像素數下，傳感器越大，單個像素都越大，入光量越多，畫質越好。

\begin{itemize}
  \item \textbf{大底感應器}：進光量大，夜景畫質好，動態範圍寬。缺點係手機要做厚（鏡頭模組凸起），成本高。
  \item \textbf{高像素小底}：白天光線充足時，高像素可以影到更清晰嘅相，方便後期裁切。但係暗光下每個像素進光不足，噪點多。
  \item \textbf{主流旗艦選擇}：主攝多數係1/1.3英寸到1英寸（部分機型已經用緊1英寸大底），像素數5000萬左右——呢個係畫質同體積嘅平衡點。
\end{itemize}

\subsubsection*{像素融合技術（Pixel Binning）：}

為咗兼顧白天高像素同夜晚大像素，手機廠商用「像素融合」——將多個小像素合併成一個大像素輸出。

\begin{itemize}
  \item \textbf{四合一像素（Quad Bayer）}：5000萬像素感應器，輸出1250萬像素相。每個輸出像素由4個相鄰像素合成，進光量提升4倍，畫質更乾淨。
  \item \textbf{九合一像素（Nonacell）}：用喺1億像素以上嘅感應器，輸出約1100-1200萬像素。同樣原理，進光量更大。
  \item \textbf{實際效果}：預設影相模式下，你攞到嘅就係融合後嘅低像素相——畫質好、檔案細。需要極致細節時（比如影文件、遠景），可以手動切到「高像素模式」輸出全解像度。
\end{itemize}

\booksubsection{八、特殊拍攝模式}

除咗日常影相，手機入面仲有幾個特殊模式，用啱可以影到平時影唔到嘅畫面。

\subsubsection*{全景模式：}

\begin{itemize}
  \item 拍攝時按提示箭頭緩慢平穩咁橫向移動手機，唔好快——快咗會出現拼接錯位。
  \item 保持手機垂直方向穩定，淨係橫向平移。如果上下晃動，地平線會變波浪。
  \item 有移動嘅人或車時慎用——拼接時會出現「半個身體」或「鬼影」。
  \item 全景模式都可以打直影——直影橫拼接，適合影高聳嘅建築或瀑布。
\end{itemize}

\subsubsection*{慢動作與延時攝影：}

\begin{itemize}
  \item \textbf{慢動作}：用120fps或240fps高幀率拍攝，回放時變成慢鏡頭。適合影水滴、運動瞬間、寵物跳躍。光線要充足，慢動作模式下進光時間短，暗光畫質差。
  \item \textbf{延時攝影}：間隔拍攝多張相合成快進影片。適合影車流、雲彩移動、日落、花朵開放。拍攝時長幾分鐘到幾個鐘不等，務必用三腳架固定手機，手持必糊。
  \item \textbf{小技巧}：延時攝影影雲彩時，將手機架好，影15-30分鐘，最後得到10秒左右嘅延時影片，效果震撼。
\end{itemize}

\subsubsection*{星空/天文模式：}

\begin{itemize}
  \item 需要遠離城市光污染，揾完全冇路燈幹擾嘅郊外或山區。城市入面影唔出銀河。
  \item 必須用三腳架，手持必糊。手機要穩穩架住幾秒到幾十秒。
  \item 天氣要晴朗無雲，月亮唔好太大（滿月會蓋住星光）。農曆月初或月末嘅暗夜最適合。
  \item iPhone嘅「夜間模式」會自動識別星空並延長曝光時間（最長30秒）；安卓旗艦嘅「星空模式」或「天文模式」效果類似。要耐心等待，撳快門後唔好鬱手機。
\end{itemize}

\subsubsection*{微距模式：}

\begin{itemize}
  \item 用於拍攝花朵、昆蟲、珠寶、文字特寫等小物體。可以貼近到1-2釐米距離對焦。
  \item 部分手機有專門嘅微距鏡頭（多數係200萬像素或500萬像素嘅小底鏡頭），畫質一般；旗艦機用超廣角鏡頭兼顧微距，畫質更好。
  \item 貼近拍攝時鏡頭離主體好近，會遮擋光線，必要時用枱燈或手電筒補光。
  \item 對焦距離極近，景深好淺，主體前後移動幾毫米就可能脱焦。拍攝時手要穩，最好等主體保持唔鬱。
\end{itemize}

\begin{tcbinfo}{靜態攝影實用建議}
\begin{itemize}
  \item 日常拍攝：HEIF/JPEG預設即可，檔案細、畫質夠用。
  \item 重要場景（風光、人像、需要後期）：開啟RAW，保留最大後期空間。
  \item 選啱鏡頭：風光用超廣角，人像用長焦，日常用主攝。
  \item 夜景同星空：務必固定手機，開夜景模式或天文模式。
  \item 大底手機（如1英寸感應器機型）畫質上限更高，但操作習慣仍係關鍵。
\end{itemize}

\end{tcbinfo}

\begin{tcbpractice}{動手試試：摸清你嘅相機能力}
\begin{enumerate}
  \item 查看你嘅手機係咪支援RAW拍攝（設定相機格式），試下開啟並拍攝一張，對比同一場景嘅JPEG同RAW檔案大小。
  \item 切換到超廣角、主攝、長焦三個焦段，分別影同一處遠景，體會視角差異同畫質差異。
  \item 揾一個有夜景模式或星空模式嘅場景，固定手機後拍攝一張夜景相。
  \item 嘗試用全景模式橫影一張大場景，注意保持手機平穩移動。
\end{enumerate}

\end{tcbpractice}
